\documentclass[11pt]{article}

\usepackage[margin=1in]{geometry}
\usepackage{amsmath,amssymb,amsfonts,bm}
\usepackage{mathtools}
\usepackage{physics}
\usepackage{enumitem}
\usepackage{booktabs}
\usepackage{hyperref}
\usepackage{xcolor}
\usepackage{algorithm}
\usepackage{algpseudocode}
\usepackage{tikz}
\usetikzlibrary{positioning,arrows.meta,fit,backgrounds,calc}

\hypersetup{
    colorlinks=true,
    linkcolor=blue,
    citecolor=blue,
    urlcolor=blue
}

% ---- notation ----
\newcommand{\CM}{\mathbf{C}_{M}}
\newcommand{\cM}{\mathbf{c}_{M}}
\newcommand{\CI}{\mathbf{C}_{I}}
\newcommand{\cI}{\mathbf{c}_{I}}
\newcommand{\NL}{\mathbf{N}^{\mathcal{L}}}
\newcommand{\rL}{\mathbf{r}^{\mathcal{L}}}
\newcommand{\BM}{\mathbf{B}_{M}}
\newcommand{\BI}{\mathbf{B}_{I}}
\newcommand{\Dd}{\mathbf{D}_{d}}
\newcommand{\Kmat}{\mathbf{K}}
\newcommand{\Rone}{\mathbf{R}_{1}}
\newcommand{\Rtwo}{\mathbf{R}_{2}}
\newcommand{\PM}{\mathbf{P}_{M}}
\newcommand{\sig}{\bm{\sigma}}
\newcommand{\eps}{\varepsilon}
\newcommand{\grad}{\nabla}
\newcommand{\Div}{\nabla\cdot}
\newcommand{\dt}{\Delta t}
\newcommand{\dg}{\Delta \gamma}
\newcommand{\Oh}{\Omega}
\newcommand{\kb}{k_{B}}
\newcommand{\Lcal}{\mathcal{L}}
\newcommand{\V}{\mathcal{V}}
\newcommand{\dV}{\,dV}
\newcommand{\dS}{\,dS}
\newcommand{\ZrMicro}{\texttt{ZrMicro}}
\newcommand{\MoDELib}{\texttt{MoDELib2-NNL}}
\newcommand{\code}[1]{\texttt{#1}}

\title{A Two-Time-Scale Coupling of the \texttt{ZrMicro} 0-D Cluster-Dynamics
Model to the \texttt{MoDELib2-NNL} Spatially-Resolved Solver for Irradiation
Growth and Creep in Zirconium}
\author{Fluor\_Zr project --- coupling formulation}
\date{\today}

\begin{document}
\maketitle

\begin{abstract}
This document specializes the two-time-scale spatial cluster-dynamics (SR-CD)
formulation to the two codes in the present project: the \textbf{0-D reduced
cluster-dynamics model implemented in \ZrMicro{}} (this repository,
\code{ZrMicro/}) and the \textbf{spatially-resolved finite-element SR-CD /
dislocation-mechanics solver \MoDELib{}}
(\url{https://github.com/mlm335/MoDELib2-NNL}), which implements the model of
Maron \emph{et al.}~\cite{Maron2024}. The mobile point-defect population
($C_v,C_i,C_{2i},C_{3i}$) is treated as a fast manifold equilibrated by the
finite-element reaction--diffusion solve in \MoDELib{}, while the immobile
cluster population is treated as a spatially varying slow manifold marched in
dose. We adopt throughout the \textbf{crystallographically resolved
representation (``Option A'')}: the interstitial $\langle a\rangle$ loops are
resolved into the three prism variants $a_1,a_2,a_3$ and the vacancy
$\langle c\rangle$ loops into one basal family, and --- \emph{for the present
stage} --- we do \textbf{not} carry \ZrMicro{}'s internal stress-aligned /
non-aligned doubling; the stress dependence enters instead through per-variant
resolved-stress weights. We (i) recast the fast/slow split in the precise state
variables and growth/creep relations coded in \ZrMicro{}; (ii) place the
construction on its footing in the numerical multiple-time-scale literature
(Tikhonov--Fenichel singular perturbation, the quasi-steady-state assumption,
computational singular perturbation, the heterogeneous multiscale method, and
projective integration); (iii) give an explicit coupling protocol with a
schematic (Fig.~\ref{fig:schematic}); (iv) document a concrete exporter,
\code{ZrMicro/py\_utils/modelib\_export.py}, that serializes a \ZrMicro{} run as
the dose-indexed closure table \MoDELib{} consumes; and (v) specify in
implementation-ready detail the \textbf{adopted operator-split QSSA scheme}
(\S\ref{sec:opsplit}) --- a steady mobile FEM solve coupled to a \emph{pointwise}
immobile dose march that reuses the \ZrMicro{} stiff integrator essentially
unchanged (one \code{freeze\_mobile} flag) and maps onto the existing OpenMP
batch mode, one case per quadrature point, needing no 0-D seed run.
This revision aligns the formulation with the D1/M1
deliverable~\cite{PoGhoniem2026D1M1}, whose \S2 and \S4 supersede the earlier
treatment on three points: the diffusion-anisotropy-difference bias is expressed
in \MoDELib{}'s $(Z^0_m,p_m)$ form, with those parameters fitted \emph{exactly}
to \ZrMicro{}'s empirical symmetric $\delta^{DAD}$ factors so the two models bias
identically (\S\ref{sec:dad}); the sub-critical vacancy
bi-pyramid enters as a third immobile sink family (\S\ref{sec:sinks},
\S\ref{sec:morph}); and vacancy-loop thermal emission is identified as a channel
the resolved model carries and the reduced model does not
(\S\ref{sec:emission}). Each is flagged where it bears on the well-mixed
regression test (\S\ref{sec:verify}).
\end{abstract}

\tableofcontents
\newpage

\section{Purpose, the Two Codes, and the Coupling at a Glance}
\label{sec:purpose}

\subsection{What each code is}

\paragraph{\ZrMicro{} (this repository --- the 0-D ``fast slow-kinetics generator'').}
\ZrMicro{} integrates a stiff system of $12$ coupled rate equations for the
volume-averaged concentrations of point defects and dislocation loops in
neutron-irradiated hcp Zr, with additional bookkeeping (six conservation
accumulators and one evolving network dislocation density), producing
irradiation growth, irradiation creep, void swelling, and radiation hardening
versus dose. It is integrated in time by SUNDIALS CVODE (BDF, dense) through a
C++ executable, with a SciPy \code{LSODA} Python fallback. It is a \emph{point}
(0-D) model: no spatial transport; grain boundaries and surfaces enter only
through lumped sink strengths.

\paragraph{\MoDELib{} (the spatial solver).}
\MoDELib{} is a C++ Mechanics-of-Defect-Evolution library whose
\code{include/ClusterDynamics/} and \code{include/FEM/} modules implement the
spatially-resolved cluster-dynamics + crystal-plasticity model of
Maron~\emph{et al.}~\cite{Maron2024}: a finite-element reaction--diffusion
boundary-value problem for the \emph{mobile} defect species on a meshed grain or
polycrystal (\code{Library/Meshes/}), with crystallographically resolved
\emph{immobile} loop families and a coupled inelastic deformation field. It
carries Zr material data in \code{Library/Materials/} and shares its elasticity,
dislocation-mechanics, and FEM infrastructure with the MoDELib2
discrete-dislocation-dynamics code base of Po and Ghoniem~\cite{Po2014MoDELib}.

\paragraph{The gap this document closes.}
\ZrMicro{} has the full nonlinear reaction kinetics and the growth/creep closure
but no spatial resolution; \MoDELib{} has the spatial FEM transport but treats
the immobile microstructure as a prescribed or saturated sink field. The
two-time-scale coupling uses \ZrMicro{}'s kinetics as the \emph{local}
slow-manifold update at each \MoDELib{} integration point, and \MoDELib{}'s FEM
solve as the \emph{fast} quasi-static mobile-species solver. This restores the
feedback loop between spatially resolved mobile transport and a spatially
evolving immobile microstructure while keeping the unknown count far below a full
spatial population balance.

\subsection{Scale separation}

The essential separation is $\tau_{M}\ll\tau_{I}$, where $\tau_M$ is the
equilibration time of the mobile species $\{v,i,2i,3i\}$ (sub-second to seconds
at reactor temperatures and dose rates) and $\tau_I$ is the dose scale over which
loop densities, radii, and the network density evolve (many dpa, i.e.\
$10^{6}$--$10^{8}\,$s at $G\sim10^{-7}$--$10^{-6}\,$dpa/s). The ratio
\begin{equation}
    \epsilon_{ts} = \frac{\tau_M}{\tau_I} \ll 1
\end{equation}
is the small parameter of the singular-perturbation construction of
\S\ref{sec:tts}. The coupled model carries two evolving manifolds,
\begin{align}
    \mathcal{M}_{M}: \quad & \CM(\mathbf{x};\mathbf{Q},\bm{\eta})
    &&\text{fast mobile concentration manifold (\MoDELib{} FEM)}, \\
    \mathcal{M}_{I}: \quad & \mathbf{Q}(\mathbf{x},t)
    &&\text{slow immobile microstructure manifold (\ZrMicro{} kinetics)},
\end{align}
with $\mathbf{Q}$ the immobile slow state (\S\ref{sec:slowstate}) and
$\bm{\eta}$ the environmental fields (temperature $T$, dose rate $G$, stress
$\sig$, grain orientation $\mathbf{R}_g$, boundary sink coefficients).

\subsection{Coupling architecture}

Figure~\ref{fig:schematic} summarizes the data flow: \ZrMicro{} runs once per
$(T,G,\sig)$ regime and exports a dose-indexed closure table; \MoDELib{} then
marches in dose, alternating a fast quasi-static FEM solve for $\CM^\star$ with a
local slow update of the crystallographically resolved loop state, and feeding
the resulting inelastic strain to its mechanics solver.

\begin{figure}[h!]
\centering
\begin{tikzpicture}[
    >=Latex, font=\small, node distance=6mm,
    box/.style={draw, rounded corners, align=center, inner sep=4pt, minimum height=8mm},
    zr/.style ={box, fill=blue!8,   draw=blue!55,   text width=33mm},
    fast/.style={box, fill=orange!12, draw=orange!70, text width=46mm},
    slow/.style={box, fill=green!10,  draw=green!55,  text width=46mm},
    mech/.style={box, fill=gray!10,   draw=gray!60,   text width=46mm},
    lbl/.style ={font=\scriptsize\itshape, align=center},
]
% ZrMicro 0-D (offline), left
\node[zr] (zr) {\textbf{\ZrMicro{} 0-D CD}\\(this repo)\\[2pt]
  \scriptsize CVODE BDF, 12 ODEs\\
  $\to \mathbf{Y}_{0D}(\gamma)$};
\node[zr, below=4mm of zr, fill=blue!4] (exp) {\scriptsize\textbf{exporter}\\
  \code{modelib\_export.py}\\
  \code{modelib\_cd\_table.txt}\\ \code{modelib\_cd\_meta.json}};
\draw[->, blue!60] (zr) -- (exp);

% MoDELib spatial loop, right column
\node[mech, right=26mm of zr, yshift=6mm] (mechB) {\textbf{Mechanics BVP} (\MoDELib{} FEM)\\
  $\Rightarrow\ \sig(\mathbf{x})$};
\node[fast, below=of mechB] (fastB) {\textbf{FAST: mobile FEM solve} (QSSA)\\
  $\mathcal{R}_M(\CM^\star;\mathbf{Q})=\mathbf{0}\Rightarrow \CM^\star(\mathbf{x})$};
\node[slow, below=of fastB] (slowB) {\textbf{SLOW: local loop update}\\
  $a_1,a_2,a_3,\langle c\rangle,\rho_N$: $\;\partial_\gamma\mathbf{Q}$
  (\ZrMicro{} kinetics)};
\node[mech, below=of slowB] (LI) {\textbf{Inelastic strain} $\mathbf{L}^{I}$
  $\Rightarrow \mathbf{F}^{I}_{n+1}$};

\draw[->] (mechB) -- (fastB);
\draw[->] (fastB) -- (slowB);
\draw[->] (slowB) -- (LI);

% feedback: next dose step
\draw[->, gray!70] (LI.east) -- ++(10mm,0) |- node[lbl, pos=0.25, right]{next dose\\$\gamma_{n+1}$} (mechB.east);

% ZrMicro feeds init/guess/calibration into the spatial loop
\draw[->, dashed, blue!60] (exp.east) -- ++(6mm,0) |-
  node[lbl, pos=0.25, below, yshift=-1mm]{$\mathbf{Y}_{0D}(\gamma)$:\\ init, Newton guess, calibration}
  (fastB.west);
\draw[->, dashed, blue!60] (exp.east) -- ++(6mm,0) |- (slowB.west);

% enclosing box for the MoDELib dose loop
\begin{scope}[on background layer]
  \node[draw=orange!40, dashed, rounded corners, fit=(mechB)(fastB)(slowB)(LI),
        inner sep=6pt, label={[lbl, orange!60!black]above:\MoDELib{} dose loop}] {};
\end{scope}
\end{tikzpicture}
\caption{Two-time-scale coupling. \ZrMicro{} (left) generates a dose-indexed
closure table via \code{modelib\_export.py}. \MoDELib{} (right) marches in dose,
alternating the \emph{fast} quasi-static mobile FEM solve (orange) with the
\emph{slow} local update of the crystallographically resolved loop state
$a_1,a_2,a_3,\langle c\rangle$ and network density $\rho_N$ (green), then
converting the microstructure rate to an inelastic strain that drives the
mechanics solve (gray). The dashed blue arrows show how the \ZrMicro{} output
initializes the slow state, seeds the Newton guess for the fast solve, and
calibrates the shared scalar kinetic coefficients.}
\label{fig:schematic}
\end{figure}

\section{The \texttt{ZrMicro} 0-D Reduced Cluster-Dynamics Model}
\label{sec:zrmicro}

This section states the \emph{actual} \ZrMicro{} model in the notation of the
code; source references are to \code{ZrMicro/py\_utils/} and
\code{ZrMicro/cpp\_utils/}.

\subsection{Mobile species (fast subset)}

The mobile defect concentration vector --- the subset that diffuses in the
spatial code --- is
\begin{equation}
\label{eq:zr_mobile}
    \CM = \begin{bmatrix} C_{v} & C_{i} & C_{2i} & C_{3i}\end{bmatrix}^{T}
    \equiv \texttt{y[0:4]},
    \qquad
    \cM = \BM \CM,
    \quad
    \BM=\mathrm{diag}(-1,1,2,3),
\end{equation}
(atom fractions). $\cM$ is the signed defect content of the mobile pool.

\subsection{Immobile species --- as coded, and the collapse used here}
\label{sec:immobile_coded}

\ZrMicro{} stores interstitial $\langle a\rangle$ loops and vacancy
$\langle c\rangle$ loops each as a \emph{non-aligned} + \emph{stress-aligned}
pair, with a number density and a defect content for every variant:
\begin{equation}
\label{eq:zr_immobile_state}
    \underbrace{C_{i\Lcal},C_{ai\Lcal},C_{v\Lcal},C_{av\Lcal}}_{\text{loop number densities}=\,\texttt{y[4:8]}},
    \qquad
    \underbrace{C_{i\Lcal}^{i},C_{ai\Lcal}^{i},C_{v\Lcal}^{v},C_{av\Lcal}^{v}}_{\text{defects in loops}=\,\texttt{y[8:12]}}.
\end{equation}
A scalar stress-alignment fraction $f_a(\sig,T)$ (\code{input\_data.py};
Woo's anisotropy~\cite{Woo1988}) partitions production between aligned and
non-aligned variants. \ZrMicro{} also evolves a network dislocation density
$\rho_N\equiv\texttt{y[18]}$ and six conservation accumulators
$\texttt{y[12:18]}$.

\paragraph{Collapse for the coupling (no aligned/non-aligned split).}
For the present stage of the coupling we do \emph{not} propagate the
aligned/non-aligned doubling into the spatial code. We collapse each pair into a
single per-species total,
\begin{equation}
\label{eq:collapse}
    N_{a\Lcal}=C_{i\Lcal}+C_{ai\Lcal},\quad
    c_{a\Lcal}=C_{i\Lcal}^{i}+C_{ai\Lcal}^{i},\quad
    N_{c\Lcal}=C_{v\Lcal}+C_{av\Lcal},\quad
    c_{c\Lcal}=C_{v\Lcal}^{v}+C_{av\Lcal}^{v},
\end{equation}
and re-resolve the interstitial total \emph{crystallographically} into the three
prism variants (\S\ref{sec:resolution}). The stress dependence that \ZrMicro{}
carried scalarly in $f_a$ is recovered, in the spatial model, through the
per-variant resolved-stress weights of \eqref{eq:variant_weight}.

\subsection{Reconstructed loop geometry}
\label{sec:radii}

Loop radii are reconstructed from content and number density
(\code{post\_process.py}):
\begin{equation}
\label{eq:zr_radii}
    r_{a\Lcal} = l_a\sqrt{\frac{c_{a\Lcal}}{N_{a\Lcal}}},
    \qquad
    r_{c\Lcal} = l_c\sqrt{\frac{c_{c\Lcal}}{N_{c\Lcal}}},
    \qquad
    l_a = \sqrt{\frac{\Oh}{\pi b_a}},
    \quad
    l_c = \sqrt{\frac{\Oh}{\pi b_c}},
\end{equation}
the discrete form of $c_I^{\Lcal}=\pi (r^{\Lcal})^2 b^{\Lcal} N^{\Lcal}$; the
line density a loop family contributes to $\rho_N$ is
$\rho^{\Lcal}=2\pi r^{\Lcal} N^{\Lcal}/\Oh$.

\subsection{Rate equations (structure)}
\label{sec:rate_structure}

Each mobile equation has the reaction form (no transport in 0-D),
\begin{equation}
\label{eq:zr_mobile_rhs}
    \frac{dC_M^{\alpha}}{dt}
    =
    G^{\alpha}
    +\tfrac12 R_2^{\alpha\beta\gamma}C_M^{\beta}C_M^{\gamma}
    + R_1^{\alpha\beta}(\mathbf{Q},\bm\eta)\,C_M^{\beta}
    + (\text{thermal emission}),
\end{equation}
with $R_2$ the binary recombination/clustering reactions, $R_1$ absorption at the
network and loops plus dissociation, and the generation vector carrying the dose
rate directly,
$\mathbf{G}=G\,[\,1-\eps_{vL},\;1-\eps_{2i}-\eps_{3i}-\eps_{iL},\;\eps_{2i},\;\eps_{3i}\,]^{T}$.
Thus $G$ [dpa/s] appears in \ZrMicro{} only through these constant rates; the
integration variable is time $t$ and dose is $\gamma=G t$
(\code{visualization.py}). The immobile equations are nucleation + flux-driven
growth + annealing; in collapsed form,
\begin{equation}
\label{eq:zr_growth_rate}
    \dot{c}_{a\Lcal}
    =
    \underbrace{\tfrac{l_c}{l}\sqrt{c_{a\Lcal}N_{a\Lcal}}\big[Z_{i}^{a}\Phi_i - g(r)Z_{v}^{a}\Phi_v\big]}_{\text{flux-driven growth}}
    + (\text{nucleation content}) + G_{a\Lcal},
\end{equation}
with mobile fluxes $\Phi_i=\omega_i(C_i+2C_{2i}+3C_{3i})$,
$\Phi_v=\omega_v C_v$, Arrhenius jump frequencies
$\omega_{i,v}=z_c\nu_{i,v}e^{-E^m_{i,v}/\kb T}$, a minimum-size gate $g(r)$, and
the diffusion-anisotropy-difference (DAD) bias factors
\begin{equation}
\label{eq:zr_DAD}
    Z_i^{a}=1+\delta^{DAD}_i,\;\;
    Z_v^{a}=1-\delta^{DAD}_v,\;\;
    Z_i^{c}=1-\delta^{DAD}_i,\;\;
    Z_v^{c}=1+\delta^{DAD}_v,
\end{equation}
($\delta^{DAD}_i$ interpolated in $T$, \code{input\_data.py}). These are the 0-D
counterparts of the orientation-resolved bias used in the spatial model
(\S\ref{sec:sinks}). Note that \eqref{eq:zr_DAD} imposes $Z^a_m+Z^c_m=2$, a
normalization the resolved bias factors do not share. In the coupled setting
$\delta^{DAD}_{i,v}$ are \emph{not} independent parameters: they are derived from
the \MoDELib{} diffusion tensors through \eqref{eq:delta_from_p}, as set out in
\S\ref{sec:dad}.

\subsection{Growth and creep closure}
\label{sec:growth_creep}

\ZrMicro{} converts loop contents to macroscopic strain (\code{post\_process.py}),
which in collapsed form reads
\begin{equation}
\label{eq:zr_growth}
    \eps_a = A_a\,c_{a\Lcal},
    \qquad
    \eps_c = -A_c\,c_{c\Lcal},
    \qquad A_a=1,\;A_c=\tfrac12,
\end{equation}
($\langle a\rangle$ loops expand basal dimensions, $\langle c\rangle$ loops
contract the $c$-axis). In the crystallographically resolved model these scalars
become a tensor source (\S\ref{sec:mechanics}); the creep response then emerges
\emph{geometrically} from the per-variant growth rates rather than from a scalar
aligned/non-aligned closure.

\section{Numerical-Methods Foundation}
\label{sec:literature}

The reduction --- freeze the fast variables on their quasi-equilibrium, march only
the slow variables --- is standard and well-analyzed. We record the lineage because
it dictates which approximations are controlled and how to verify them.

\paragraph{Singular perturbation / geometric reduction.}
In standard slow--fast form
$\epsilon_{ts}\dot{\CM}=\mathcal F_M(\CM,\mathbf Q)$,
$\dot{\mathbf Q}=\mathcal F_I(\CM,\mathbf Q)$,
Tikhonov's theorem~\cite{Tikhonov1952} guarantees that after a fast transient
solutions track $\mathcal F_M=\mathbf0$ when that manifold is asymptotically
stable in the fast variables; Fenichel's geometric singular perturbation
theory~\cite{Fenichel1979} makes the slow manifold
$\mathcal M_M=\{\mathcal F_M=\mathbf0\}$ normally hyperbolic and persistent for
small $\epsilon_{ts}>0$, justifying the $\epsilon_{ts}\to0$ limit of
\S\ref{sec:tts}. The relevant fast Jacobian is the mobile reaction Jacobian
\eqref{eq:reaction_jacobian}; its eigenvalues set $\tau_M$ and must stay stable.

\paragraph{Quasi-steady-state assumption (QSSA).}
Setting $\mathcal F_M=\mathbf0$ is the QSSA. Segel and
Slemrod's analysis~\cite{SegelSlemrod1989} is the canonical demonstration that
the QSSA is a singular-perturbation limit with a quantifiable small parameter;
the present model is the QSSA applied to the mobile defects, which lets us treat
$\CM^\star$ as an elliptic-BVP function of the slow state.

\paragraph{ILDM and computational singular perturbation.}
When the split is not obvious it can be computed: the intrinsic low-dimensional
manifold method of Maas and Pope~\cite{MaasPope1992} and the computational
singular perturbation method of Lam and Goussis~\cite{LamGoussis1994} identify
and refine slow/fast subspaces from the Jacobian spectrum. Here the split is
physical (mobile vs.\ immobile); CSP/ILDM are the tools to \emph{verify} that no
immobile mode becomes fast (or mobile mode slow) at extreme $T,G,\sig$.

\paragraph{Heterogeneous multiscale and equation-free methods.}
The macro/micro pattern --- a coarse spatial solver calling a fine local solver for
closure on demand --- is the heterogeneous multiscale method of E and
Engquist~\cite{EEngquist2003} and the equation-free framework of Kevrekidis
\emph{et al.}~\cite{Kevrekidis2003}. Our coupling is an HMM whose ``micro'' solver
is the local \ZrMicro{} kinetics and whose ``macro'' solver is the \MoDELib{} FEM
transport plus mechanics. Projective / coarse integration~\cite{GearKevrekidis2003}
is the formal justification for taking the dose step $\dg$ much larger than
$\tau_M$.

\paragraph{Operator splitting, IMEX, reduced-order acceleration.}
Advancing the stiff fast residual separately from the slow update is operator
splitting~\cite{Strang1968}; treating the stiff reaction implicitly and
slow/transport terms explicitly is the IMEX strategy of Ascher, Ruuth, and
Spiteri~\cite{AscherRuuthSpiteri1997} --- \ZrMicro{} already uses an implicit stiff
integrator (CVODE BDF) for exactly this reason. Because the fast BVP is solved
repeatedly, proper generalized decomposition / reduced-basis
surrogates~\cite{Chinesta2011} are a natural accelerator (\S\ref{sec:improvements}).

\section{Two-Time-Scale Reformulation}
\label{sec:tts}

\subsection{Fast and slow split}

Promoting \ZrMicro{}'s 0-D kinetics to a field theory,
\begin{align}
\label{eq:full_fast_slow_1}
    \epsilon_{ts}\,\partial_t \CM
    &=
    \mathcal{F}_{M}(\CM,\mathbf{Q};\bm{\eta})
    = -\Div\mathbf J_M + \tfrac12 \Rtwo[\CM,\CM] + \Rone(\mathbf Q,\bm\eta)\CM + \PM,
    \\
\label{eq:full_fast_slow_2}
    \partial_t \mathbf{Q}
    &=
    \mathcal{F}_{I}(\CM,\mathbf{Q};\bm{\eta}),
\end{align}
where \eqref{eq:full_fast_slow_1} differs from the 0-D
equation~\eqref{eq:zr_mobile_rhs} only by the transport term $-\Div\mathbf J_M$.
Taking $\epsilon_{ts}\to0$,
\begin{equation}
\label{eq:fast_constraint}
    \mathcal{F}_{M}(\CM^\star,\mathbf{Q};\bm{\eta}) = \mathbf{0}
    \quad\text{(fast quasi-static BVP, \MoDELib{} FEM)},
    \qquad
    \partial_t \mathbf{Q}
    =
    \mathcal{F}_{I}(\CM^{\star},\mathbf{Q};\bm{\eta})
    \;\;\text{(slow, \ZrMicro{})}.
\end{equation}

\subsection{The crystallographically resolved slow state}
\label{sec:slowstate}

We adopt the slow state with the three prism $\langle a\rangle$ interstitial-loop
variants and one basal $\langle c\rangle$ vacancy-loop family per grain,
\begin{equation}
\label{eq:Q_state}
    \mathbf{Q}(\mathbf x,\gamma)
    =
    \big\{
    \underbrace{N_{a_1\Lcal},N_{a_2\Lcal},N_{a_3\Lcal},N_{c\Lcal}}_{\text{number densities}},\;
    \underbrace{c^{a_1}_I,c^{a_2}_I,c^{a_3}_I,c^{c}_I}_{\text{contents}},\;
    \rho_N,\;
    \underbrace{n_v}_{\text{morphology}}
    \big\},
\end{equation}
with habit normals $\hat{\mathbf a}_1,\hat{\mathbf a}_2,\hat{\mathbf a}_3$ (the
three $\langle 11\bar20\rangle$ a-directions at $120^\circ$ in the basal plane)
and $\hat{\mathbf c}$. No aligned/non-aligned doubling is carried;
\eqref{eq:Q_state} has the same number of loop families as \ZrMicro{}
\eqref{eq:zr_immobile_state} but indexes them by crystallography rather than by
stress alignment.

\subsection{Fast mobile manifold (the \MoDELib{} solve)}

For fixed $\mathbf Q$ the per-species strong residual is
\begin{equation}
\label{eq:mobile_residual}
    \mathcal{R}_{M}^{\alpha}
    =
    -\Div\mathbf{J}_{M}^{\alpha}
    + \tfrac{1}{2}R_{2}^{\alpha\beta\gamma}C_{M}^{\beta}C_{M}^{\gamma}
    + R_{1}^{\alpha\beta}(\mathbf{Q},\bm{\eta})C_{M}^{\beta}
    + P_{M}^{\alpha}(\bm{\eta}),
    \quad \alpha\in\{v,i,2i,3i\},
\end{equation}
with the anisotropic hcp flux
\begin{equation}
\label{eq:mobile_flux}
    \mathbf{J}_{M}^{\alpha}
    =
    -\mathbf{D}^{\alpha}(\mathbf x)\grad C_{M}^{\alpha}
    -\frac{\mathbf{D}^{\alpha}C_{M}^{\alpha}}{\kb T}\grad\Psi^{\alpha},
    \quad
    \mathbf{D}^{\alpha}
    =
    \mathbf{R}_{g}\,\mathrm{diag}(D_a^\alpha,D_a^\alpha,D_c^\alpha)\,\mathbf{R}_{g}^{T},
\end{equation}
$\mathbf R_g(\mathbf x)$ the grain orientation, $\Psi^\alpha$ an optional drift
potential. The first-order operator $R_1(\mathbf Q,\bm\eta)$ is built from the
\emph{local} loop sink strengths (\S\ref{sec:sinks}), which is what makes the fast
solve see a spatially varying microstructure.

\subsection{Slow immobile manifold (the local \ZrMicro{} update)}
\label{sec:slow_eqs}

Per loop family $X\in\{a_1,a_2,a_3,c\}$,
\begin{equation}
\label{eq:slow_content_component}
    \frac{\partial c_{I}^{X}}{\partial t}
    =
    \underbrace{\sum_{m}\langle k^2_{X,m}\rangle D_m B_M^{m} C_M^{m,\star}}_{\text{absorption of local }\CM^\star}
    + P_I^{X} + S_I^{X,\mathrm{nuc}} - S_I^{X,\mathrm{diss}},
\end{equation}
the field form of \ZrMicro{}'s growth equation \eqref{eq:zr_growth_rate} with the
local $\CM^\star$ in place of the 0-D mean. Number densities evolve by the same
nucleation/annealing laws,
\begin{equation}
\label{eq:nucleation_law}
    \frac{\partial N^{X}}{\partial t}
    =
    \mathcal{J}^{X}_{\mathrm{nuc}}(\CM^{\star},T,\sig,G)
    - \mathcal{K}^{X}_{\mathrm{anneal}}(N^{X},T)
    - \mathcal{K}^{X}_{\mathrm{coal}}(N^{X},r^{X})
    - \mathcal{K}^{X}_{\mathrm{abs}}(N^{X},\Gamma),
\end{equation}
where the first three terms are \ZrMicro{}'s nucleation, annealing, and
coalescence, and $\mathcal K_{\mathrm{abs}}$ is a new spatial sink at grain
boundaries / interfaces $\Gamma$ (zero recovers \ZrMicro{}). The radius rate
follows from $c_I^{X}=\pi(r^{X})^2 b^{X}N^{X}$,
\begin{equation}
\label{eq:r_evolution_general}
    \frac{\partial r^{X}}{\partial t}
    =
    \frac{1}{2\pi r^{X}b^{X}N^{X}}
    \left[\frac{\partial c_{I}^{X}}{\partial t} - \pi (r^{X})^{2}b^{X}\frac{\partial N^{X}}{\partial t}\right].
\end{equation}

\subsection{Dose as the slow march variable}

Evolve the slow manifold in dose,
\begin{equation}
\label{eq:dose_slow_update}
    d\gamma = G\,dt,
    \qquad
    \frac{\partial \mathbf{Q}}{\partial \gamma}
    = \frac{1}{G}\,\mathcal{F}_{I}(\CM^{\star},\mathbf{Q};\bm{\eta}),
\end{equation}
matching \ZrMicro{}'s dose-indexed histories and the exporter table of
\S\ref{sec:exporter}.

\section{Crystallographic Resolution of the Interstitial Loops (Option A)}
\label{sec:resolution}

This is the single modeling decision of the coupling: how the collapsed
interstitial total $(N_{a\Lcal},c_{a\Lcal})$ of \eqref{eq:collapse} is resolved
into the three prism variants.

\subsection{Resolved-stress variant weights}

Production into variant $k$ is weighted by the Boltzmann factor of the resolved
normal stress $\sigma_{kk}=\hat{\mathbf a}_k\cdot\sig\cdot\hat{\mathbf a}_k$ on
its habit plane,
\begin{equation}
\label{eq:variant_weight}
    w_k(\sig,T)
    =
    \frac{\exp(\sigma_{kk}\Oh/\kb T)}{\sum_{k'}\exp(\sigma_{k'k'}\Oh/\kb T)},
    \qquad \sum_{k=1}^{3} w_k = 1,
\end{equation}
which reduces to the isotropic $w_k=\tfrac13$ at zero deviatoric stress. This
replaces \ZrMicro{}'s scalar alignment fraction $f_a$ \emph{and removes the need
for the aligned/non-aligned doubling}: stress anisotropy is expressed directly as
an imbalance among $w_1,w_2,w_3$.

\paragraph{Relation to the SIPN partition of the resolved model.}
Equation~\eqref{eq:variant_weight} takes the activation volume to be the scalar
atomic volume $\Oh$. The resolved model instead uses the full stress-induced
preferential-nucleation (SIPN) partition, Eq.~(37)
of~\cite{PoGhoniem2026D1M1}, in which the cascade-produced defects are
distributed over the immobile families through the \emph{formation-volume
tensor} $\bm\Omega_k$ of each family,
\begin{equation}
\label{eq:sipn}
    \chi^k_{i}
    =
    \frac{\exp\!\big(\sig\!:\!\bm\Omega_k/\kb T\big)\,I^k_{i}}
         {\sum_{p}\exp\!\big(\sig\!:\!\bm\Omega_p/\kb T\big)\,I^p_{i}},
    \qquad
    \chi^k_{v}
    =
    \frac{\exp\!\big(\sig\!:\!\bm\Omega_k/\kb T\big)\,I^k_{v}}
         {\sum_{p}\exp\!\big(\sig\!:\!\bm\Omega_p/\kb T\big)\,I^p_{v}},
\end{equation}
where $I^k_i=\tfrac12(1+p_k)$ and $I^k_v=\tfrac12(1-p_k)$ are the interstitial and
vacancy selectors built from the family polarity $p_k$ ($p=+1$ for the three
$\langle a\rangle$ variants, $p=-1$ for the $\langle c\rangle$ family), so that
the two partitions in \eqref{eq:sipn} normalize independently over the
interstitial and vacancy families. For a pure prismatic loop $\bm\Omega_k$
reduces to
$\bm\Omega_k=(\mathbf b_k\otimes\mathbf b_k/\mathbf b_k\!\cdot\!\mathbf b_k)V_k$,
so $\sig\!:\!\bm\Omega_k=\sigma_{kk}V_k$ and \eqref{eq:sipn} collapses onto
\eqref{eq:variant_weight} with the activation volume identified as the cluster
volume $V_k=m_k\Oh_0$ rather than $\Oh$. The two therefore agree in form but not
in scale, and \eqref{eq:variant_weight} should be read as the single-defect
($m_k=1$) case. The linearized weights
$w_{a_k}(\sig)=\tfrac13\big(1+\chi\hat\sigma_{kk}\big)$ of Eq.~(36)
of~\cite{PoGhoniem2026D1M1} are in turn the small-stress limit of
\eqref{eq:sipn}; all three forms coincide at $\sig=\mathbf0$, giving the common
seed $\tfrac13$. Because the pyramid limit
$\bm\Omega_k=(\mathbf 1/3)V_k$, with $\mathbf 1$ the second-order identity, is
\emph{spherical}, $\sig\!:\!\bm\Omega_k=-q\,V_k$ then depends only on the
hydrostatic part of $\sig$: sub-critical vacancy embryos are
partitioned by pressure alone, and acquire a deviatoric preference only as they
flatten onto the basal plane (\S\ref{sec:morph}).

\subsection{Initialization and the lumped recovery}

At initialization (and in the offline table of \S\ref{sec:exporter}) the per-variant
fields are
\begin{equation}
\label{eq:variant_split}
    N_{a_k\Lcal} = w_k\,N_{a\Lcal},
    \qquad
    c^{a_k}_I = w_k\,c_{a\Lcal},
    \qquad
    r_{a_k\Lcal} = l_a\sqrt{\frac{c^{a_k}_I}{N_{a_k\Lcal}}} = l_a\sqrt{\frac{c_{a\Lcal}}{N_{a\Lcal}}},
\end{equation}
so the three variants share the 0-D mean radius (the weight cancels in the
radius) until the spatial kinetics differentiate them. Summing the variants
recovers the \ZrMicro{} lumped interstitial loop population,
$\sum_k N_{a_k\Lcal}=N_{a\Lcal}$ and $\sum_k c^{a_k}_I=c_{a\Lcal}$ --- the basis of
the well-mixed regression test (\S\ref{sec:verify}). The vacancy
$\langle c\rangle$ loops are a single family, $N_{c\Lcal}=N_{c\Lcal}$,
$c^{c}_I=c_{c\Lcal}$, with morphology $n_v=|c^{c}_I|/N_{c\Lcal}$ (vacancies per
loop).

\subsection{Once resolved, the variants evolve independently}

After initialization, \MoDELib{} integrates each variant with
\eqref{eq:slow_content_component}--\eqref{eq:r_evolution_general} using its
\emph{own} habit-plane orientation: the bias factors entering
$\langle k^2_{a_k,m}\rangle$ depend on the angle between $\hat{\mathbf a}_k$ and
the local diffusion-anisotropy / stress axes, so the variants diverge in density
and radius wherever the stress or texture is non-trivial. The shared scalar
coefficients ($\mathcal J_{\mathrm{nuc}}$, $\tau$, $\delta^{DAD}$, coalescence)
are calibrated once from \ZrMicro{} and reused across all three variants.

\subsection{Diffusion-anisotropy consistency}
\label{sec:dad}

The scalar DAD biases $\delta^{DAD}_{i,v}(T)$ \eqref{eq:zr_DAD} and the tensor
diffusivities $D^\alpha_{a,c}(T)$ \eqref{eq:mobile_flux} describe the same
physics, but they are \emph{not} two parameterizations of one identity. In
\MoDELib{} the bias is not a parameter at all: it is generated by the
anisotropy of the species diffusion tensor through
\begin{equation}
\label{eq:pm_dad}
    p_m=\Big(\tfrac{D_c^m}{D_a^m}\Big)^{1/6},
    \quad
    Z_{a\Lcal,m}^{\mathrm{DAD},0}=Z_m^0\,\tfrac{p_m+p_m^{-2}}{2},
    \quad
    Z_{c\Lcal,m}^{\mathrm{DAD},0}=Z_m^0\,p_m,
\end{equation}
which is Eq.~(15) of the D1/M1 deliverable~\cite{PoGhoniem2026D1M1}. \ZrMicro{}
instead \emph{posits} $\delta^{DAD}_{i}(T),\delta^{DAD}_{v}$ phenomenologically.
Section~4.3 of~\cite{PoGhoniem2026D1M1} is explicit that the two coincide only in
the isotropic limit $p_m\to1$, $\delta\to0$, and that away from it the map
$\delta\leftrightarrow(p,Z^0)$ ``is a modelling choice and must be interpreted
rather than taken as an identity.''

\paragraph{Direction of calibration (adopted).}
We take the \textbf{0-D empirical forms as the reference} and fit the resolved
model's $(Z^0_m,p_m)$ to reproduce them, rather than deriving $\delta^{DAD}$ from
the migration tensors. The information therefore flows
$\delta^{DAD}\!\to\!(Z^0_m,p_m)$, and \eqref{eq:pm_dad} is used as a
\emph{parameterization} of the bias rather than as a prediction from the
diffusivities.

The fit is \emph{exact}, not a least-squares compromise: \eqref{eq:pm_dad}
carries two free parameters per species and \eqref{eq:zr_DAD} imposes two targets
$(Z^a_m,Z^c_m)$, so the system inverts in closed form. Eliminating $Z^0_m$
between the two relations of \eqref{eq:pm_dad} gives $2Z^a_m/Z^c_m=1+p_m^{-3}$,
whence
\begin{equation}
\label{eq:delta_from_p}
    p_m=\Big(\frac{2Z^a_m}{Z^c_m}-1\Big)^{-1/3},
    \qquad
    Z^0_m=\frac{Z^c_m}{p_m},
\end{equation}
evaluated with the 0-D targets $Z^a_i=1+\delta^{DAD}_i$, $Z^c_i=1-\delta^{DAD}_i$
for the interstitial species and $Z^a_v=1-\delta^{DAD}_v$,
$Z^c_v=1+\delta^{DAD}_v$ for the vacancy. At the reference condition
$T=573$~K, where $\delta^{DAD}_i=0.072113$ and $\delta^{DAD}_v=0.107886$, this
yields
\begin{equation}
\label{eq:dad_fitted}
    p_v=1.178808,\;\; Z^0_v=0.939836,
    \qquad
    p_i=p_{2i}=p_{3i}=0.913720,\;\; Z^0_i=1.015504,
\end{equation}
reproducing the four 0-D factors to machine precision ($\sim\!10^{-16}$). The
fitted anisotropies are moreover physically sensible for h.c.p. Zr: $p_v>1$
places the faster vacancy diffusion along $c$, and $p_i<1$ places the faster
interstitial diffusion in the basal plane.

\paragraph{What the fit costs.}
Two properties of \eqref{eq:delta_from_p} must be recorded, because they are the
price of exact agreement.

First, \textbf{$p_m$ is no longer the diffusional anisotropy of the material
file.} In \code{Library/Materials/Zr4.txt} the migration-energy tensors are
\emph{isotropic for both monomers} --- $E^m_v=0.9$~eV and $E^m_i=0.09$~eV on all
three diagonal components, with isotropic $D_0$ --- and anisotropic only for the
di- and tri-interstitials ($0.041$~eV in the basal plane against $0.187$~eV along
$c$), so the tensors would give
\begin{equation}
\label{eq:pm_values}
    p_v=p_i=1,
    \qquad
    p_{2i}=p_{3i}=\exp\!\Big[-\frac{E^{m}_{33}-E^{m}_{11}}{6\kb T}\Big]\simeq0.61
    \quad(T=573\;\mathrm{K}).
\end{equation}
The fitted values \eqref{eq:dad_fitted} differ from these throughout --- most
starkly for the monomers, where the tensors imply \emph{no} DAD bias at all
against fitted $p_v=1.179$, $p_i=0.914$. The bias in the coupled model is thus
\textbf{calibrated to the reduced model, not generated from the tensorial flux}.
The physical claim of \S2 of~\cite{PoGhoniem2026D1M1} --- that the resolved
formulation obtains DAD ``naturally'' from the anisotropic diffusion --- is
therefore \emph{not} what the coupled runs exercise; recovering it would require
migration tensors consistent with \eqref{eq:dad_fitted}, which is a separate
materials-data question.

Second, the calibration is \textbf{tied to one temperature}. The 0-D
$\delta^{DAD}_i(T)$ is interpolated linearly between $300$~K and $600$~K anchors
\eqref{eq:zr_DAD}, so \eqref{eq:dad_fitted} holds at $573$~K only; $p_m$ and
$Z^0_m$ must be refitted whenever the run temperature changes. The underlying
tension is the normalization: \eqref{eq:zr_DAD} forces $Z^a_m+Z^c_m=2$, which
\eqref{eq:pm_dad} does not satisfy for any $p_m$ unless $Z^0_m$ absorbs the
difference --- which is exactly what \eqref{eq:delta_from_p} makes it do.

\paragraph{Status.}
With \eqref{eq:dad_fitted} in \code{Zr4.txt} the two models now bias the same
channels by the same factors, so the well-mixed regression of
\S\ref{sec:verify} is meaningful \emph{as it stands} --- no 0-D refit is
required, and a residual discrepancy in $\eps_a,\eps_c$ between the resolved
interior and the 0-D reference is now a genuine finding rather than expected
bookkeeping.

\subsection{Sink strengths from the local slow state}
\label{sec:sinks}

The first-order operator is assembled locally from $\mathbf Q(\mathbf x)$:
\begin{equation}
\label{eq:R1_local}
    \Rone = \Rone\big[\{N_{a_k\Lcal},r_{a_k\Lcal}\}_{k=1}^3, N_{c\Lcal},r_{c\Lcal},\rho_N,T,\sig\big],
\end{equation}
with loop and network sink strengths
\begin{equation}
\label{eq:loop_sinks}
    \langle k^2_{X,m}\rangle
    =
    Z_{X,m}^{\mathrm{tot}}(\sig,T,\mathbf D^m)\,\underbrace{2\pi r_{X}N_{X}}_{S^{\mathrm{geo}}_X},
    \qquad
    \langle k^2_{N,m}\rangle = Z_{N,m}\,\rho_N,
\end{equation}
and the total bias the DAD factor times a stress-induced preferential-absorption
(SIPA) increment,
$Z_{X,m}^{\mathrm{tot}}=Z_{X,m}^{\mathrm{DAD},0}\big[1+\Delta Z^{\mathrm{SIPA}}_{X,m}(\sig)/Z_{X,m}^{\mathrm{DAD},0}\big]$.
Each $\langle a\rangle$ variant evaluates its SIPA increment at its \emph{own}
resolved stress $\sigma_{kk}$ and the $\langle c\rangle$ family at the basal
value, which is Eqs.~(46)--(49) of~\cite{PoGhoniem2026D1M1}; all four reduce to
their stress-free DAD values \eqref{eq:pm_dad} at $\sig=\mathbf0$. The geometric
factor $S^{\mathrm{geo}}_X=2\pi r_X N_X$ is exactly what the exporter table
tabulates (\S\ref{sec:exporter}); \MoDELib{} multiplies it by the local
$Z_{X,m}^{\mathrm{tot}}$.

\paragraph{The bi-pyramid family in the sink sum.}
The total sink strength seen by mobile species $m$ runs over \emph{three}
immobile families, not two: Eq.~(11) of~\cite{PoGhoniem2026D1M1} sums over
$\mathcal I\!\Lcal\in\{c\Lcal,a\Lcal,bp\}$, where the compact sub-critical
vacancy bi-pyramids contribute
\begin{equation}
\label{eq:k2_bp}
    \langle k^2_{bp,m}\rangle=\alpha_{bp}\,s_{bp}\,\pi r_{bp}N_{bp},
\end{equation}
a volumetric (three-dimensional) capture law, in contrast to the planar
$2\pi r N$ scaling of \eqref{eq:loop_sinks}. This family has no counterpart in
\ZrMicro{}, whose $\langle c\rangle$ population is a single loop family, and it
is \emph{not} carried by the present exporter table (\S\ref{sec:exporter}). Two
consequences follow. First, the closure table as currently written under-reports
the vacancy sink strength wherever the $\langle c\rangle$ population is still
sub-critical --- i.e. at low dose, exactly where the denuded-zone structure is
established. Second, the interpolation of \S\ref{sec:morph} is what connects
\eqref{eq:k2_bp} to \eqref{eq:loop_sinks}: $\langle k^2_{bp,m}\rangle$ and
$\langle k^2_{c\Lcal,m}\rangle$ are the two limits between which the vacancy
sink strength is blended as the embryos flatten.

\subsection{Vacancy bi-pyramid$\to$loop morphology}
\label{sec:morph}

The $\langle c\rangle$-cluster morphology variable is the mean number of
vacancies per cluster, $\bar m_{c}=c^{c}_I/N_{c\Lcal}$ (written $n_v$ in the
exporter table), and a transition $S(\bar m_c)\in[0,1]$ interpolates every
$\langle c\rangle$ property $Y$ between bi-pyramid and basal-loop limits,
\begin{equation}
\label{eq:interp}
    Y_{cl}=S(\bar m_c)\,Y_{c\Lcal}+[1-S(\bar m_c)]\,Y_{bp},
    \qquad
    S(\bar m_c)=\frac{1}{1+\exp[-(\bar m_c-m_0)/w]},
\end{equation}
including sink strength, relaxation volume, and eigenstrain. This is the
transition Maron \emph{et al.}~\cite{Maron2024} identify as necessary for
breakaway growth; the spatial code carries it per integration point because
$\bar m_c$ is now a field.

\paragraph{Calibration of the transition (as implemented in the resolved model).}
Eqs.~(28)--(33) of~\cite{PoGhoniem2026D1M1} fix both the argument and the
endpoints of \eqref{eq:interp}, and the coupling must adopt them rather than
treat $a_s,n_s$ as free. The midpoint and width are set by the stability window
$[m_{\min},m_{\max}]$,
\begin{equation}
\label{eq:sigmoid_cal}
    m_0=\frac{m_{\min}+m_{\max}}{2}-m_s,
    \qquad
    w=w_0\big(m_{\max}-m_{\min}\big),
\end{equation}
so that clusters with $\bar m_c\lesssim m_{\min}$ are predominantly
bi-pyramidal and those with $\bar m_c\gtrsim m_{\max}$ predominantly planar. The
two geometric endpoints are
\begin{equation}
\label{eq:radii_limits}
    R_{bp}=\Big(\frac{\bar m_c\Oh_0}{\sqrt8}\Big)^{1/3},
    \qquad
    R_{c\Lcal}=\sqrt{\frac{\bar m_c\Oh_0}{\pi b_c}},
\end{equation}
i.e. a volumetric $\bar m_c^{1/3}$ law crossing over to the planar
$\bar m_c^{1/2}$ law of \eqref{eq:zr_radii}. The \emph{same} sigmoid
interpolates the formation-volume tensor between its spherical and planar limits,
\begin{equation}
\label{eq:omega_interp}
    \bm\Omega_k=\big[1-S(\bar m_c)\big]\frac{\mathbf 1}{3}V_k
    + S(\bar m_c)\,\frac{\mathbf b_k\otimes\mathbf b_k}{\mathbf b_k\!\cdot\!\mathbf b_k}V_k,
    \qquad V_k=\bar m_c\Oh_0,
\end{equation}
which is what couples the morphology back into the SIPN partition
\eqref{eq:sipn}: a sub-critical embryo responds only to pressure, and acquires
its deviatoric (basal-normal) preference continuously as $S\to1$. Consistently,
the vacancy sink strength blends \eqref{eq:k2_bp} into \eqref{eq:loop_sinks}
through the same $S$. A free-energy form,
$S(\bar m_c,T,\sig)=\{1+\exp[(G_{c\Lcal}-G_{bp})/\kb T]\}^{-1}$, lets the
transition shift with local stress, temperature, and chemistry.

\subsection{Vacancy-loop thermal emission: a channel the 0-D model omits}
\label{sec:emission}

The resolved model removes vacancies from the $\langle c\rangle$ family through
\emph{two} distinct channels, and the distinction matters for the coupling
because \ZrMicro{} implements only one of them. The first is the dissolution of
sub-critical embryos that fail to reach the minimum stable size, at the Arrhenius
lifetime $\tau_{c\Lcal}(T)=\tau^0\exp(E_a/\kb T)$; this removes both density and
content and is the term \ZrMicro{} carries. The second, Eqs.~(54)--(57)
of~\cite{PoGhoniem2026D1M1}, is the \emph{thermal emission} of single vacancies
from an already-stable loop, evaluated at the mean size,
\begin{equation}
\label{eq:loop_emission}
    \alpha_{v,c\Lcal}(\bar m_c)
    =
    \beta_{v,c\Lcal}\exp\!\Big[-\frac{E^{b}_{v}(\bar m_c)}{\kb T}\Big],
\end{equation}
with $E^{b}_{v}(\bar m_c)$ the size-dependent vacancy binding energy, which grows
toward its bulk limit so that larger loops emit ever more weakly. This channel
reduces $c^{c}_I$ \emph{without} removing loops, so it enters the content
equation but leaves the density equation untouched --- unlike the dissolution
term, which debits both. The vacancies released by both channels are returned to
the mobile pool through the source $\mathbf S_{v\Lcal}$, so the exchange is a
pure transfer and the point-defect budget stays closed: a surviving loop may
simultaneously grow by absorbing the excess vacancy flux and shrink by thermal
emission without double counting.

\ZrMicro{} has no analogue of \eqref{eq:loop_emission}: its loop binding energy
is a single constant $\approx5$~eV surrogate (Table~1
of~\cite{PoGhoniem2026D1M1}, ``pending a refit''), which is size-independent and
therefore cannot reproduce the weakening emission of a growing loop. Two
implications for the coupling. First, the well-mixed regression of
\S\ref{sec:verify} compares a resolved model that emits against a reduced model
that does not, so agreement is expected to degrade at high temperature and high
dose, where $\bar m_c$ is large and \eqref{eq:loop_emission} is most active ---
in the direction of \MoDELib{} predicting \emph{smaller} $\langle c\rangle$ loops
than \ZrMicro{}. Second, the material file must supply the binding energies that
\eqref{eq:loop_emission} and the mobile-species emission matrix require; the key
\code{mobileSpeciesBindingEnergy\_eV} is referenced by Eq.~(9)
of~\cite{PoGhoniem2026D1M1} but is absent from the current
\code{Library/Materials/Zr4.txt}.

\section{Weak Form and Newton Solve of the Fast Manifold}
\label{sec:weak}

\MoDELib{}'s FEM module assembles the weak form of \eqref{eq:mobile_residual}.
With test function $\widehat C^\alpha$,
\begin{equation}
\label{eq:weak_fast}
    \int_{\V}\grad\widehat C^{\alpha}\cdot\mathbf J_M^{\alpha}\dV
    + \int_{\V}\widehat C^{\alpha}\Big[\tfrac12 R_2^{\alpha\beta\gamma}C_M^\beta C_M^\gamma + R_1^{\alpha\beta}C_M^\beta + P_M^\alpha\Big]\dV
    - \int_{\partial\V}\widehat C^{\alpha}\,\mathbf J_M^{\alpha}\cdot\mathbf n\dS = 0.
\end{equation}

\subsection{Boundary conditions}

Use the Robin (finite-rate sink) condition
\begin{equation}
\label{eq:robin_bc}
    -\mathbf J_M^\alpha\cdot\mathbf n = h_\alpha\,(C_M^\alpha - C_{\mathrm{eq}}^\alpha)
    \quad\text{on }\Gamma_h,
\end{equation}
which recovers the perfect-sink Dirichlet condition as $h_\alpha\to\infty$ and
no-flux as $h_\alpha\to0$. \ZrMicro{} corresponds to $h_\alpha\to\infty$ lumped
into the sinks. A thermodynamically consistent reservoir concentration follows
from the defect chemical potential
$\mu^\alpha=\mu_0^\alpha+\kb T\ln C_M^\alpha-\sig:\Delta\bm\Omega^\alpha$,
\begin{equation}
\label{eq:ceq_mu}
    C_{\mathrm{eq}}^\alpha
    = \exp\!\Big[\frac{\mu_\Gamma^\alpha-\mu_0^\alpha+\sig:\Delta\bm\Omega^\alpha}{\kb T}\Big],
\end{equation}
coupling stress through the formation-volume tensor rather than normal traction
alone.

\subsection{Newton linearization}

The fast BVP is nonlinear through $R_2$. At iteration $j$, solve
$\mathbf K_{MM}(\CM_j,\mathbf Q)\,\delta\CM = -\mathbf R_M(\CM_j,\mathbf Q)$ with
the reaction Jacobian (symmetric $R_2$)
\begin{equation}
\label{eq:reaction_jacobian}
    \frac{\partial\mathcal R_M^\alpha}{\partial C_M^\delta}
    = -\Div\Big(\frac{\partial\mathbf J_M^\alpha}{\partial C_M^\delta}\Big)
    + R_2^{\alpha\delta\gamma}C_M^\gamma + R_1^{\alpha\delta}.
\end{equation}
This is the fast-subsystem Jacobian whose spectrum sets $\tau_M$
(\S\ref{sec:literature}); monitoring its conditioning checks QSSA validity. For
positivity, solve for $u^\alpha=\ln C_M^\alpha$ or apply a positivity line
search; clip only as a last resort and log it.

\section{Mechanics Coupling}
\label{sec:mechanics}

The inelastic velocity gradient decomposes as
$\mathbf L^I=\mathbf L^V+\mathbf L^D+\mathbf L^P$. With the loops resolved by
crystallography, the plastic part is a clean tensor sum over the four habit
normals,
\begin{equation}
\label{eq:LP_Zr}
    \mathbf L^P
    =
    \dot c_I^{a_1}(\hat{\mathbf a}_1\otimes\hat{\mathbf a}_1)
    +\dot c_I^{a_2}(\hat{\mathbf a}_2\otimes\hat{\mathbf a}_2)
    +\dot c_I^{a_3}(\hat{\mathbf a}_3\otimes\hat{\mathbf a}_3)
    +\dot c_I^{c}(\hat{\mathbf c}\otimes\hat{\mathbf c}).
\end{equation}
Volume-averaging \eqref{eq:LP_Zr} over a stress-free grain (where
$\dot c_I^{a_k}=\tfrac13\dot c_{a\Lcal}$) reproduces the \ZrMicro{} scalar growth
\eqref{eq:zr_growth}: $\tfrac13(\hat{\mathbf a}_1\otimes\hat{\mathbf a}_1+\dots)
=\tfrac12(\mathbf I-\hat{\mathbf c}\otimes\hat{\mathbf c})$ gives the basal
$\eps_a$ and $\hat{\mathbf c}\otimes\hat{\mathbf c}$ gives $\eps_c$. Under stress,
the variant imbalance \eqref{eq:variant_weight} makes \eqref{eq:LP_Zr}
deviatoric, producing irradiation creep \emph{geometrically} --- no separate creep
closure is needed. The volumetric and diffusive parts are
\begin{equation}
\label{eq:LVD}
    \mathbf L^V
    = \sum_n\frac{\Delta\Oh_n^0}{3\Oh_0}\dot c_M^n\,\mathbf I
    + \sum_{X}f^{X}\partial_t\big[(\mathbf b^{X}\otimes\mathbf A^{X})N^{X}\big],
    \qquad
    J^I\,\mathrm{tr}(\mathbf L^D)=\sum_n\Div\mathbf J_M^n.
\end{equation}
Swelling and growth rates are the spherical and deviatoric parts of $\mathbf L^I$.
The stress returned by the \MoDELib{} mechanics solve re-enters the fast manifold
through the BC \eqref{eq:ceq_mu}, the bias \eqref{eq:loop_sinks}, the variant
weights \eqref{eq:variant_weight}, and the morphology transition --- closing the
loop.

\section{Coupling Protocol: \texttt{ZrMicro} $\leftrightarrow$ \texttt{MoDELib2-NNL}}
\label{sec:coupling}

\subsection{Roles and data exchange}

\begin{table}[h!]
\centering
\small
\begin{tabular}{@{}p{0.28\linewidth} p{0.33\linewidth} p{0.31\linewidth}@{}}
\toprule
\textbf{Task} & \textbf{\ZrMicro{} (0-D, this repo)} & \textbf{\MoDELib{} (spatial FEM)} \\
\midrule
Fast mobile solve $\CM^\star(\mathbf x)$ & provides $R_1,R_2,\PM$ and a 0-D Newton guess & \code{include/ClusterDynamics/}+\code{include/FEM/}: weak form \eqref{eq:weak_fast}, Newton \eqref{eq:reaction_jacobian} \\
Local slow update $\partial_\gamma\mathbf Q$ & \code{rate\_equations} RHS \eqref{eq:slow_content_component}--\eqref{eq:nucleation_law} as per-point micro-model & integrates $a_1,a_2,a_3,c$ per integ.\ point \\
Growth/creep $\to\mathbf L^I$ & \code{post\_process} closure \eqref{eq:zr_growth} & tensor source \eqref{eq:LP_Zr}, mechanics BVP \\
Init \& calibration & dose table $\mathbf Y_{0D}(\gamma)$ via \code{modelib\_export.py} & seeds $\mathbf Q(\mathbf x,\gamma_0)$ and closure coefficients \\
Material data & \code{input/*.xlsx} (3 sheets) & \code{Library/Materials/} (Zr), \code{Library/Meshes/} \\
\bottomrule
\end{tabular}
\caption{Division of labor in the two-time-scale coupling (cf.\ Fig.~\ref{fig:schematic}).}
\end{table}

\subsection{The exporter: \texttt{modelib\_export.py}}
\label{sec:exporter}

The ``Level 1'' deliverable is \code{ZrMicro/py\_utils/modelib\_export.py}, which
serializes a completed \ZrMicro{} run as the closure table \MoDELib{} reads. It
performs exactly the collapse \eqref{eq:collapse} and the resolved-stress split
\eqref{eq:variant_split}, and writes two files:
\begin{itemize}[leftmargin=1.4em,topsep=2pt,itemsep=1pt]
\item \code{modelib\_cd\_table.txt} --- a whitespace-delimited, \code{\#}-commented
table with one row per dose point and columns
\[
\big[\,\gamma,t,\;C_v,C_i,C_{2i},C_{3i},\;
\{N_{a_k},N_c\},\{c^{a_k}_I,c^{c}_I\},\{r_{a_k},r_c\},
\{S^{\mathrm{geo}}_{a_k},S^{\mathrm{geo}}_c\},\;\rho_N,n_v,\eps_a,\eps_c\,\big],
\]
number densities and contents in $\mathrm{m^{-3}}$ (atom fraction $/\Oh$), radii
in m, $S^{\mathrm{geo}}_X=2\pi r_X N_X$ in $\mathrm{m^{-2}}$;
\item \code{modelib\_cd\_meta.json} --- scalars \MoDELib{} needs to rebuild $R_1$
and $\mathbf L^I$: $T,G,\Oh,b_a,b_c,l_a,l_c$, the activation volume, the DAD bias
factors $Z_i^a,Z_v^a,Z_i^c,Z_v^c,Z_N$, the variant weights $w_k$ and the resolved
stress used, the habit normals $\hat{\mathbf a}_k,\hat{\mathbf c}$, the column
schema, and units.
\end{itemize}
Both are plain text, parsed trivially by \MoDELib{}'s C++ I/O, numpy, or pandas.
Usage from the \ZrMicro{} notebook after \code{process\_solution}:
\begin{verbatim}
    from py_utils.modelib_export import export_for_modelib
    export_for_modelib(results, input_data, run_dir)                 # zero stress
    export_for_modelib(results, input_data, run_dir,                 # uniaxial:
                       sigma_a_resolved=(s1, s2, s3))                #  resolved
                                                                     #  on a_k
\end{verbatim}
At zero stress the weights are $w_k=\tfrac13$; under a resolved stress they tilt
toward the more-stressed variant via \eqref{eq:variant_weight}.

\paragraph{Known gaps in the current schema.}
The table above was fixed before~\cite{PoGhoniem2026D1M1} and does not yet carry
everything \S\ref{sec:sinks}--\S\ref{sec:emission} require. Three columns/blocks
are missing, and each is a silent under-report rather than an error the reader
will notice:
\begin{enumerate}[label=(\roman*),leftmargin=1.6em,itemsep=1pt]
\item \textbf{The bi-pyramid family.} No $N_{bp},r_{bp}$ or
$S^{\mathrm{geo}}_{bp}$, so \eqref{eq:k2_bp} cannot be assembled and the vacancy
sink strength is under-reported at low dose.
\item \textbf{The morphology state.} $n_v$ is exported, but not $S(\bar m_c)$ nor
the calibration $(m_{\min},m_{\max},m_s,w_0)$ of \eqref{eq:sigmoid_cal}, so
\MoDELib{} cannot reproduce the blend that \eqref{eq:interp} and
\eqref{eq:omega_interp} depend on.
\item \textbf{The DAD scalars are the superseded ones.} The exported
$Z_i^a,Z_v^a,Z_i^c,Z_v^c$ are \ZrMicro{}'s phenomenological values, which
\S\ref{sec:dad} shows are not realizable from the resolved diffusion tensors.
Under the adopted calibration direction these should not be exported as inputs to
\MoDELib{} at all --- \MoDELib{} generates its own bias from $\mathbf D^m$ ---
and are better retained as \emph{diagnostics} recording what the reduced model
assumed.
\end{enumerate}
Until the schema is extended, the table remains valid for $\mathbf Q$
initialization and for the geometric factors $S^{\mathrm{geo}}_X$, which is what
the operator-split scheme of \S\ref{sec:opsplit} actually consumes.

\subsection{Three levels of integration}

\begin{enumerate}[label=(\arabic*),leftmargin=1.6em]
\item \textbf{Offline table (implemented).} \ZrMicro{} emits
\code{modelib\_cd\_table.txt} per $(T,G,\sig)$ regime; \MoDELib{} reads it as a
look-up closure for $\mathbf Q$ initialization and for the shared scalar kinetic
coefficients, keyed on the local $(T,G,\sig)$.
\item \textbf{File/array interface.} The same tables drive both initialization and
the Newton initial guess $\CM^{(0)}=\CM^{0D}(\gamma)$ via the spatial correction
$\CM=\CM^{0D}+\delta\CM$, with $\mathcal R_M(\CM^{0D}+\delta\CM;\mathbf Q)=\mathbf0$.
\item \textbf{On-line HMM call.} For tightest coupling, replace the look-up by a
direct call to \ZrMicro{}'s RHS evaluated with the local $\CM^\star$: expose
\code{rate\_equations.\_rhs\_full} (Python) or the C++ \code{integrate\_one}
(\code{rate\_equations.cpp}) as the per-point micro-model. The C++ path is
preferred and already builds a standalone executable with a \code{--key=value}
CLI and an OpenMP \code{--batch\_file} mode that can evaluate many integration
points concurrently in one subprocess.
\end{enumerate}

\subsection{Dose-stepping algorithm}
\label{sec:algorithm}

\begin{algorithm}[h!]
\caption{Two-time-scale coupled SR-CD march (\ZrMicro{} kinetics, \MoDELib{} FEM)}
\begin{algorithmic}[1]
\State Run \ZrMicro{} per regime; export $\mathbf Y_{0D}(\gamma)$ with \code{modelib\_export.py}.
\State Initialize $\mathbf Q_0(\mathbf x)$ from the table: split $\langle a\rangle$ into $a_1,a_2,a_3$ by $w_k(\sig_0)$ \eqref{eq:variant_split}; set $N_c,c^c_I,\rho_N,n_v$.
\State Initialize mechanics fields $\mathbf F^I_0,\mathbf u_0,\sig_0$.
\For{$n=0,1,\dots,N-1$}
    \State $\dg=\gamma_{n+1}-\gamma_n$;\; $\dt=\dg/G_n$.
    \State Solve the \MoDELib{} mechanics BVP with $\mathbf F^I_n$ to get $\sig_n(\mathbf x)$.
    \State At each integ.\ point update $\mathbf D^\alpha_n,\mathbf Z_n,\Kmat_n,\Rone_n,R_{2,n}$ from $\mathbf Q_n,T_n,G_n,\sig_n$ (\S\ref{sec:resolution}, \S\ref{sec:sinks}); refresh $w_k(\sig_n)$.
    \State Build BCs: Robin \eqref{eq:robin_bc} with $C^\alpha_{\mathrm{eq}}$ from \eqref{eq:ceq_mu}.
    \State \textbf{Fast solve:} $\CM^{(0)}\leftarrow\CM^{0D}(\gamma_n)$; Newton-iterate \eqref{eq:reaction_jacobian} with positivity line search until $\|\mathbf R_M\|<\mathrm{tol}_M$, giving $\CM_n^\star(\mathbf x)$.
    \State \textbf{Slow update:} for each $X\in\{a_1,a_2,a_3,c\}$ evaluate $\dot c^{X}_I,\dot N^{X}$ from \eqref{eq:slow_content_component}--\eqref{eq:nucleation_law} and $\dot r^{X}$ from \eqref{eq:r_evolution_general}; update $\rho_N$.
    \State March $\mathbf Q_{n+1}=\mathbf Q_n+\dt\,\dot{\mathbf Q}_n$ (explicit), or projective integration~\cite{GearKevrekidis2003} for large $\dg$.
    \State Update $n_v,S(n_v)$ and the interpolated $\langle c\rangle$ properties \eqref{eq:interp}.
    \State Form $\mathbf L^P_{n+1}$ \eqref{eq:LP_Zr}, $\mathbf L^V_{n+1},\mathbf L^D_{n+1},\mathbf L^I_{n+1}$; update $\mathbf F^I_{n+1}=\exp(\dt\,\mathbf L^I_{n+1})\mathbf F^I_n$.
    \State Output $\CM_n^\star,\mathbf Q_{n+1}$, sink strengths, growth/swelling strain, diagnostics.
\EndFor
\end{algorithmic}
\end{algorithm}

\section{Adopted Implementation: Operator-Split QSSA
(Steady Mobile $\oplus$ Pointwise Immobile)}
\label{sec:opsplit}

This is the implementation the project will build. It is a Lie--Trotter operator
split~\cite{Strang1968} of the two-time-scale system \eqref{eq:full_fast_slow_1}--\eqref{eq:full_fast_slow_2}:
the fast (mobile) block is solved to steady state on the FE mesh (spatially
coupled), and the slow (immobile) block is advanced \emph{pointwise} with the
\emph{existing} \ZrMicro{} stiff integrator, holding the mobile species frozen at
their local steady-state values. It requires no 0-D seed run and reuses the
\ZrMicro{} reaction kinetics and CVODE machinery essentially unchanged.

\subsection{Structure of the split}
\label{sec:opsplit_structure}

Over one dose step $[\gamma_n,\gamma_{n+1}]$, $\dg=\gamma_{n+1}-\gamma_n$,
$\dt=\dg/G$, perform a frozen-coefficient sweep:
\begin{enumerate}[label=\textbf{S\arabic*.},leftmargin=2.2em,itemsep=2pt]
\item \textbf{Mechanics} (constant here, kept general): solve the \MoDELib{}
mechanics BVP with $\mathbf F^I_n$ to obtain $\sig_n(\mathbf x)$.
\item \textbf{Assemble local coefficients} from the current immobile state
$\mathbf Q_n(\mathbf x)$ and $\sig_n$: diffusivities $\mathbf D^\alpha$, biases
$\mathbf Z$, the first-order sink operator $\Rone(\mathbf Q_n)$
\eqref{eq:R1_local}, the quadratic $\Rtwo$, the variant weights $w_k(\sig_n)$
\eqref{eq:variant_weight}, and the boundary data \eqref{eq:robin_bc}.
\item \textbf{Fast mobile solve} (FEM, Newton; spatially coupled): with
$\mathbf Q_n$ \emph{fixed}, solve $\mathcal R_M(\CM^\star;\mathbf Q_n)=\mathbf 0$
\eqref{eq:weak_fast},\eqref{eq:reaction_jacobian} for $\CM^\star_n(\mathbf x)$.
\item \textbf{Slow immobile march} (pointwise; the existing \ZrMicro{} integrator):
at every quadrature point $q$ integrate $d\mathbf Q/dt=\mathbf g(\mathbf Q;\CM^\star_n(\mathbf x_q),\bm\eta_q)$
\eqref{eq:reduced_rhs} over $[t_n,t_n+\dt]$ with $\CM^\star$ \emph{frozen},
giving $\mathbf Q_{n+1}(\mathbf x_q)$.
\item \textbf{Post-process per point}: radii \eqref{eq:zr_radii}, $n_v$,
$S(n_v)$, and the inelastic source $\mathbf L^I_{n+1}$ \eqref{eq:LP_Zr},
\eqref{eq:LVD}; update $\mathbf F^I_{n+1}=\exp(\dt\,\mathbf L^I_{n+1})\mathbf F^I_n$.
\end{enumerate}
Steps S2--S3 couple space (through diffusion); step S4 does not (loops do not
diffuse), which is the structural fact exploited in \S\ref{sec:opsplit_batch}.

\subsection{The frozen-mobile reduced right-hand side}
\label{sec:opsplit_rhs}

At a quadrature point, the slow block is the \ZrMicro{} immobile RHS with the
mobile concentrations replaced by the local steady value $\CM^\star$. For each
loop family $X\in\{a_1,a_2,a_3,c\}$,
\begin{equation}
\label{eq:reduced_rhs}
    \mathbf g:\quad
    \frac{dN^{X}}{dt}
    = \underbrace{\mathcal J^{X}_{\mathrm{nuc}}(\CM^\star,\sig,G)}_{\text{nucleation, incl.\ }w_k}
      - \mathcal K^{X}_{\mathrm{anneal}}(N^{X},T)
      - \mathcal K^{X}_{\mathrm{coal}}(\mathbf Q),
    \qquad
    \frac{dc_{I}^{X}}{dt}
    = \dot c^{X,\mathrm{abs}}_I(\CM^\star) + S^{X,\mathrm{nuc}}_I(\CM^\star) - S^{X,\mathrm{diss}}_I,
\end{equation}
together with the network density $d\rho_N/dt=\mathcal C_{\mathrm{LN}}(\mathbf Q)$
(loop--network coalescence). The absorption source is the \emph{same} sink flux
that the mobile solve balanced,
\begin{equation}
\label{eq:abs_consistent}
    \dot c^{X,\mathrm{abs}}_I
    = \sum_{m\in\{v,i,2i,3i\}}\langle k^2_{X,m}\rangle\,D_m\,B_M^{m}\,C_M^{m,\star},
\end{equation}
with $\langle k^2_{X,m}\rangle=Z^{\mathrm{tot}}_{X,m}\,2\pi r_X N_X$
\eqref{eq:loop_sinks}; the growth fluxes $\Phi_i=\omega_i(C_i^\star+2C_{2i}^\star+3C_{3i}^\star)$,
$\Phi_v=\omega_v C_v^\star$ and the homogeneous nucleation rates
($\propto R_{i,3i},R_{2i,2i}$, weighted by $w_k$ into the three prism variants)
are all evaluated at the frozen $\CM^\star$.

\paragraph{Implementation trick (minimal code change).}
Do \emph{not} write a new RHS. Integrate the full \ZrMicro{} state vector
$\texttt{y[0:19]}$ as today, but gate the first four entries with a boolean
\code{freeze\_mobile}: when set, force
\[
    \dot{\texttt{y}}[0]=\dot{\texttt{y}}[1]=\dot{\texttt{y}}[2]=\dot{\texttt{y}}[3]=0
\]
in \code{rate\_equations.cpp::rhs} (and the Python \code{\_rhs\_full}) \emph{after}
all reaction terms are evaluated. Initialize $\texttt{y[0:4]}=\CM^\star$. The
mobile slots then remain exactly $\CM^\star$ for the whole sub-integration, while
the immobile equations $\texttt{y[4:12]}$, the conservation accumulators
$\texttt{y[12:18]}$, and $\rho_N=\texttt{y[18]}$ read them as frozen parameters.
This is the entire change needed to turn the 0-D solver into the local slow
solver: one flag and a four-line zeroing. The immobile block is also far less
stiff than the full system (the stiffest modes were mobile recombination), so
CVODE takes large internal steps and the per-point solve is cheap.

\subsection{Pointwise decoupling and batch-mode mapping}
\label{sec:opsplit_batch}

Once $\CM^\star(\mathbf x)$ is known, the immobile ODEs at different quadrature
points are \emph{independent}. Step S4 therefore maps one-to-one onto the
existing OpenMP \code{--batch\_file} mode of \code{solver.exe}: \textbf{one case
per quadrature point}. The driver (\code{cpp\_bridge.run\_cpp\_solver\_batch})
writes one line per point with the tokens
\begin{verbatim}
   freeze_mobile=1  t_begin=<t_n>  t_end=<t_n+dt>  n_points=2
   y0_0=<Cv*>  y0_1=<Ci*>  y0_2=<C2i*>  y0_3=<C3i*>        # frozen mobile
   y0_4=<N_a1> ... y0_11=<c_c>  y0_18=<rho_N>              # immobile state Q_n(x_q)
   T=<T>  sigma=<sigma_n(x_q)>  G=<G>  w1=<w1> w2=<w2> w3=<w3>
\end{verbatim}
and reads back the endpoint $\mathbf Q_{n+1}(\mathbf x_q)$ (the last output row)
per case, in input order. All quadrature points integrate concurrently in a
single subprocess, each on its own \code{SUNContext}/vectors/integrator memory,
exactly as the temperature batch does today. No SUNDIALS objects are shared
across threads.

\subsection{Algorithm}
\label{sec:opsplit_algo}

\begin{algorithm}[h!]
\caption{Operator-split QSSA march (steady mobile FEM $\oplus$ pointwise immobile CVODE)}
\begin{algorithmic}[1]
\State \textbf{Initialize} $\mathbf Q_0(\mathbf x)$ pristine (\S\ref{sec:opsplit_init}); $\mathbf F^I_0=\mathbf I$, $\sig_0$.
\State Choose dose grid $\{\gamma_n\}$ with small $\dg$ through the early transient (\S\ref{sec:opsplit_dt}).
\For{$n=0,1,\dots,N-1$}
    \State $\dg=\gamma_{n+1}-\gamma_n$;\; $\dt=\dg/G$.
    \State \textbf{[S1]} Solve mechanics with $\mathbf F^I_n\Rightarrow\sig_n(\mathbf x)$.
    \State \textbf{[S2]} At each integ.\ point assemble $\mathbf D^\alpha,\mathbf Z,\Rone(\mathbf Q_n),\Rtwo,w_k(\sig_n)$ and BCs.
    \State \textbf{[S3]} \emph{Fast solve:} $\CM^{(0)}\!\leftarrow\!\CM^\star_{n-1}$ (or local reaction balance);
    \Statex \hspace{2.2em} Newton-iterate \eqref{eq:reaction_jacobian} with positivity line search until $\|\mathbf R_M\|<\mathrm{tol}_M\Rightarrow\CM^\star_n(\mathbf x)$.
    \State \textbf{[S4]} \emph{Slow march (batch):} for all points $q$ in parallel, integrate \eqref{eq:reduced_rhs}
    \Statex \hspace{2.2em} with \code{freeze\_mobile=1}, $\texttt{y}[0\!:\!4]=\CM^\star_n(\mathbf x_q)$, over $[t_n,t_n+\dt]\Rightarrow\mathbf Q_{n+1}(\mathbf x_q)$.
    \State \textbf{[split control]} If $\max_q \|\mathbf Q_{n+1}-\mathbf Q_n\|/\|\mathbf Q_n\|>\eta_{\mathrm{split}}$: halve $\dg$, redo from S3 (\S\ref{sec:opsplit_dt}).
    \State \textbf{[S5]} Per point: radii, $n_v$, $S(n_v)$; form $\mathbf L^I_{n+1}$ \eqref{eq:LP_Zr},\eqref{eq:LVD}; $\mathbf F^I_{n+1}=\exp(\dt\,\mathbf L^I_{n+1})\mathbf F^I_n$.
    \State \textbf{Output} $\CM^\star_n,\mathbf Q_{n+1}$, sink strengths, $\eps_a,\eps_c$, conservation residual (\S\ref{sec:opsplit_cons}).
\EndFor
\end{algorithmic}
\end{algorithm}

\subsection{Splitting error and dose-step control}
\label{sec:opsplit_dt}

Freezing $\CM^\star$ over the step ignores the drift of $\mathbf Q$ (hence of the
sinks) within $[\,t_n,t_n+\dt\,]$; this is a first-order Lie--Trotter splitting
error, $\mathcal O(\dg)$ per step, one-sided (with strengthening sinks the loops
are slightly over-grown). Control it by the slow state, \emph{not} by CVODE's
internal step:
\begin{itemize}[leftmargin=1.4em,topsep=2pt,itemsep=1pt]
\item \textbf{Step cap.} Accept the step only if the relative change in the sink
strength is small, e.g.\ $\max_q\,\Delta\langle k^2\rangle/\langle k^2\rangle\le\eta_{\mathrm{split}}$
with $\eta_{\mathrm{split}}\sim0.05$--$0.1$; otherwise halve $\dg$ and redo from S3.
The outer $\dg$ is bounded only by this and by how fast $\sig$ moves --- the
immobile stiffness is absorbed inside CVODE.
\item \textbf{Strang (second order), optional.} Replace S3--S4 by
[half immobile step $\dt/2$] $\to$ [re-solve mobile] $\to$ [half immobile step
$\dt/2$], giving $\mathcal O(\dg^2)$ at the cost of one extra fast solve.
\item \textbf{In-step Picard (fully coupled step), optional.} Iterate S3$\leftrightarrow$S4
to convergence of $(\CM^\star,\mathbf Q)$ before advancing; removes the splitting
error entirely and is the natural fallback if a regime ever needs it.
\end{itemize}
For the first implementation use plain Lie--Trotter with the step cap.

\subsection{Conservation across the split}
\label{sec:opsplit_cons}

In steady state the mobile pool does not accumulate, so production balances
recombination, loop absorption, and \emph{boundary outflux}:
\begin{equation}
\label{eq:cons_split}
    \int_{\V} G^{\mathrm{tot}}\dV
    =
    \int_{\V}\!\Big(\mathcal R_{\mathrm{rec}} + \sum_X \dot c^{X,\mathrm{abs}}_I\Big)\dV
    + \underbrace{\oint_{\partial\V}\sum_m \mathbf J_M^m\cdot\mathbf n\dS}_{\text{boundary loss (new channel)}}.
\end{equation}
Because the immobile absorption \eqref{eq:abs_consistent} uses the same
$\langle k^2_{X,m}\rangle D_m C_M^{m,\star}$ that the mobile solve balanced, atoms
are conserved to splitting order \emph{provided} the volume check
\eqref{eq:cons_split} includes the boundary integral. The 0-D model folded that
flux into $\rho_N$; here it is explicit. Reuse the \ZrMicro{} accumulators
$\texttt{y[12:18]}$ for the per-point volumetric terms and add the FE boundary
integral of $\mathbf J_M\cdot\mathbf n$.

\subsection{Initialization without the 0-D code}
\label{sec:opsplit_init}

Start pristine at $\gamma_0\approx0$ (after cascade thermalization), letting the
slow march build the microstructure itself:
\begin{equation}
\label{eq:pristine}
    N^{X}_0 = N_{\mathrm{embryo}}\;(\text{small}),
    \quad
    c^{X}_{I,0}\to 0^{+}\;(\Rightarrow r^X=r_{\min}),
    \quad
    \rho_{N,0}=\rho_N(T)\;(\text{grown-in}),
    \quad
    n_{v,0}=n_{\min},
\end{equation}
with accumulated strain $\eps_a=\eps_c=0$ (path-consistent, since contents start
at zero). No 0-D run is needed: the reduced RHS \eqref{eq:reduced_rhs} \emph{is}
the \ZrMicro{} nucleation+growth kinetics, so it regenerates the loop population.
The price is small $\dg$ through the early transient (where, in the 0-D code,
CVODE took its smallest steps); the step cap of \S\ref{sec:opsplit_dt} sizes this
automatically. An optional single 0-D run remains useful only to provide the
well-mixed regression reference (\S\ref{sec:verify}) and a starting $\gamma_0$
large enough to skip the transient.

\subsection{Per-quadrature-point storage and reuse checklist}
\label{sec:opsplit_storage}

\paragraph{Persistent (carried across dose steps), per quadrature point:}
the immobile slow state $\mathbf Q=\{N^{a_1},N^{a_2},N^{a_3},N^{c},c^{a_1}_I,c^{a_2}_I,c^{a_3}_I,c^{c}_I,\rho_N\}$,
the inelastic deformation $\mathbf F^I$, and the previous mobile solution
$\CM^\star_{n-1}$ (Newton warm start).

\paragraph{Transient (recomputed each step):}
$\CM^\star$, fluxes $\mathbf J_M^\alpha$, radii $r^X$, $n_v$, $S(n_v)$, sink
strengths $\langle k^2_{X,m}\rangle$, and $\mathbf L^P,\mathbf L^V,\mathbf L^D,\mathbf L^I$.

\paragraph{Reuse checklist.}
\begin{enumerate}[label=(\alph*),leftmargin=1.6em,itemsep=1pt]
\item Add \code{freeze\_mobile} to \code{parameters.h}, \code{rate\_equations.cpp::rhs},
and \code{py\_utils/rate\_equations.py::\_rhs\_full} (four-line zeroing of the
mobile derivatives, \S\ref{sec:opsplit_rhs}).
\item Keep \emph{all} immobile terms when frozen --- growth and nucleation (use
$\CM^\star$) \emph{and} the $\CM$-independent annealing $\tau(T)$, coalescence,
and $\rho_N$ evolution. Do not silently drop the latter.
\item Drive step S4 through \code{cpp\_bridge.run\_cpp\_solver\_batch} with one
case per quadrature point (\S\ref{sec:opsplit_batch}); request only the endpoint
(\code{n\_points=2}, read the last row).
\item Resolve $\langle a\rangle$ into $a_1,a_2,a_3$ with $w_k(\sig)$
\eqref{eq:variant_split} at initialization; thereafter evolve the three variants
independently (their bias factors differ by habit-plane orientation).
\item Newton warm start $=\CM^\star_{n-1}$; first step, the diffusion-free local
reaction balance.
\item Keep the existing \code{CvL\_v}/\code{CavL\_v} slight-negative clamp inside
the immobile integration.
\item Conservation check per \eqref{eq:cons_split}, including the boundary flux.
\end{enumerate}

\section{Verification and Diagnostics}
\label{sec:verify}

\paragraph{Recovery of the 0-D model (primary regression test).}
Set $h_\alpha\to\infty$ lumped into the sinks, $\mathcal K_{\mathrm{abs}}=0$,
$w_k=\tfrac13$ (zero stress), and uniform fields with periodic BCs. The fast BVP
then has no transport gradient and \eqref{eq:full_fast_slow_1} collapses to
\eqref{eq:zr_mobile_rhs}; summing the three prism variants
($\sum_k N_{a_k\Lcal}=N_{a\Lcal}$, $\sum_k c^{a_k}_I=c_{a\Lcal}$) the spatial
march should reproduce the \ZrMicro{} mobile concentrations, loop densities,
radii, and growth $\eps_a,\eps_c$. The exporter writes $\eps_a,\eps_c$ precisely
so this comparison is one column subtraction.

\paragraph{What this test can and cannot establish.}
The transport, reaction, production, and emission operators map exactly under the
volume-average identification, so in the well-mixed limit they agree to
integrator tolerance and any discrepancy there is a genuine coupling defect.
Three ingredients, however, do \emph{not} map exactly, and each admits a nonzero
residual that is physics rather than bug:
\begin{enumerate}[leftmargin=1.6em]
\item \textbf{The DAD bias} (\S\ref{sec:dad}). \emph{No longer a source of
residual}: $(Z^0_m,p_m)$ are fitted through \eqref{eq:delta_from_p} to reproduce
\eqref{eq:zr_DAD} exactly, so both models apply identical capture efficiencies to
identical channels. A discrepancy here now indicates a real defect. Note however
that the fit is tied to $T=573$~K and must be redone at other temperatures.
\item \textbf{The bi-pyramid family} (\eqref{eq:k2_bp}). Absent from the reduced
model and from the exporter, it adds vacancy sink strength at low dose that
\ZrMicro{} does not carry.
\item \textbf{Vacancy-loop thermal emission} (\S\ref{sec:emission}). Present only
in the resolved model, biting hardest at high $T$ and high dose.
\end{enumerate}
Accordingly the regression should be run \emph{first} at low temperature, low
dose, and small $\bar m_c$, where (ii) and (iii) are quiescent and the test is a
clean check of the coupling machinery, and only then extended. Running it first
at the $573$~K / high-dose condition of the current closure table conflates all
three effects with any actual defect.

\paragraph{Per-step monitors.}
\begin{align}
    \text{mobile residual / Jacobian cond.:}\;&\|\mathcal R_M\|,\ \kappa(\mathbf K_{MM})
        &&\text{(QSSA validity, \S\ref{sec:literature})}\\
    \text{positivity:}\;&\min_{\mathbf x,\alpha}C_M^\alpha(\mathbf x)>0,\\
    \text{slow-state bounds:}\;&\min N^{X}\ge0,\ \min r^{X}>0,\ 0\le S\le1,\ \textstyle\sum_k w_k=1,\\
    \text{sign/mass consistency:}\;&\cM=\BM\CM,\ \cI=\BI\CI\ \text{(\ZrMicro{} accumulators \texttt{y[12:18]})},\\
    \text{mechanics consistency:}\;&\mathbf L^I=\mathbf L^P+\mathbf L^V+\mathbf L^D.
\end{align}

\paragraph{Staged activation.}
After the well-mixed regression passes, switch on capabilities one at a time:
(i) finite-rate boundary sinks $h_\alpha<\infty$; (ii) spatial gradients in
$T,G,\sig$; (iii) resolved-stress variant imbalance $w_k\neq\tfrac13$ and the
geometric creep response \eqref{eq:LP_Zr}; (iv) local nucleation and variable
$N^{X}$; (v) free-energy morphology transition; (vi) reduced-order acceleration.
Each stage is validated against the previous in its common limit.

\section{Suggested Refinements}
\label{sec:improvements}

\begin{itemize}[leftmargin=1.4em]
\item \textbf{Reconcile the fitted $(Z^0_m,p_m)$ with the migration tensors.}
\eqref{eq:dad_fitted} reproduces the 0-D bias exactly but is inconsistent with
the anisotropy the material file's $E^m_{ij}$ imply \eqref{eq:pm_values}, so the
coupled runs do not exercise the ``DAD generated from the tensorial flux''
mechanism of~\cite{PoGhoniem2026D1M1}. Either refit the migration tensors to
\eqref{eq:dad_fitted}, or accept $(Z^0_m,p_m)$ as an independent
parameterization and say so.
\item \textbf{Make the DAD fit temperature-aware.} \eqref{eq:dad_fitted} holds at
$573$~K only, because $\delta^{DAD}_i(T)$ is $T$-interpolated in \ZrMicro{};
evaluate \eqref{eq:delta_from_p} at run temperature instead of hard-coding it.
\item \textbf{Carry the bi-pyramid family through the exporter}
\eqref{eq:k2_bp} together with the sigmoid calibration
\eqref{eq:sigmoid_cal}, so the resolved model can assemble the vacancy sink
strength it expects.
\item \textbf{Add size-dependent loop binding energies} to replace the constant
$5$~eV surrogate, enabling \eqref{eq:loop_emission} in the reduced model and
supplying \code{mobileSpeciesBindingEnergy\_eV} to the material file.
\item \textbf{Finite-rate boundary sinks} \eqref{eq:robin_bc} for imperfect grain
boundaries, oxide barriers, and solute-decorated interfaces.
\item \textbf{Chemical-potential stress coupling} \eqref{eq:ceq_mu} with the
formation-volume tensor.
\item \textbf{Per-variant radii from kinetics.} The shared-radius initialization
\eqref{eq:variant_split} is only a starting condition; once stress differentiates
the variants, let each $r_{a_k}$ evolve independently via
\eqref{eq:r_evolution_general}.
\item \textbf{Boundary-layer mesh refinement} where the depletion length
$\ell_m\sim(\langle k^2_{\Lcal,m}\rangle+\langle k^2_{N,m}\rangle)^{-1/2}\ll L$.
\item \textbf{PGD / reduced basis}~\cite{Chinesta2011} for the repeated fast
solves, optionally with $T,G,k^2$ as parametric coordinates.
\item \textbf{Re-introduce the aligned/non-aligned split later} if a regime
demands it: it becomes a second weighting layer \emph{within} each prism variant
and does not alter the structure above.
\item \textbf{Verify the fast/slow split} with a CSP/ILDM
check~\cite{LamGoussis1994,MaasPope1992} at the $T,G$ extremes.
\end{itemize}

\section{Summary}

The coupling reuses each code for what it does well: \ZrMicro{} supplies the full
nonlinear cluster-dynamics kinetics and the growth/creep closure as a
\emph{local} micro-model, and \MoDELib{} supplies the finite-element transport and
dislocation mechanics as the \emph{spatial} macro-solver. Cast as a
two-time-scale / heterogeneous-multiscale problem, the mobile point defects form
a fast quasi-static manifold solved by FEM at each dose checkpoint, while the
immobile loop population --- resolved into the three prism $\langle a\rangle$
variants and the basal $\langle c\rangle$ family (Option A), without the
aligned/non-aligned doubling --- forms a spatially varying slow manifold marched in
dose. Stress anisotropy enters through resolved-stress variant weights, so
irradiation creep emerges geometrically from the per-variant growth rather than
from a scalar closure. The single implemented bridge is
\code{modelib\_export.py}, which serializes a \ZrMicro{} run as the dose-indexed
table \MoDELib{} consumes. The adopted implementation (\S\ref{sec:opsplit}) is an
operator-split QSSA: a steady, spatially-coupled mobile FEM solve alternating
with a \emph{pointwise} immobile dose march that reuses the \ZrMicro{} stiff
integrator unchanged (gated by a single \code{freeze\_mobile} flag) and runs one
case per quadrature point through the existing OpenMP batch mode; it needs no 0-D
seed --- a pristine start regenerates the microstructure --- and its only
approximation is a first-order splitting error controlled by the dose-step cap.
The construction reduces exactly to the validated 0-D model in the well-mixed
limit, which is the primary regression test, and is grounded in the standard
numerical theory of singular-perturbation / QSSA reductions, HMM, and projective
integration.

\begin{thebibliography}{99}

% ---- materials / project references ----
\bibitem{Maron2024}
Maron, M., Li, Y., Lalani, I., Baker, K., Ramirez Flores, B., Black, T.,
Hollenbeck, J., Ghoniem, N., and Po, G. (2024).
\newblock Spatially-resolved cluster dynamics modeling of irradiation growth in
zirconium.
\newblock \emph{International Journal of Plasticity}, \textbf{175}, 103943.

\bibitem{PoGhoniem2026D1M1}
Po, G. and Ghoniem, N. (2026).
\newblock Deliverable D1/M1: Simulation of Initial Irradiation Damage.
\newblock Multiscale Materials Solutions (MMS), 29 July 2026.
\newblock \code{Docs/Formulation/GW\_Phase4\_D1M1\_NG.pdf}.

\bibitem{LiGhoniem2020}
Li, Y. and Ghoniem, N. (2020).
\newblock Cluster dynamics modeling of irradiation growth in single crystal Zr.
\newblock \emph{Journal of Nuclear Materials}, \textbf{540}, 152312.

\bibitem{Po2022}
Po, G., Huang, Y., Li, Y., Baker, K., Flores, B. R., Black, T., Hollenbeck, J.,
and Ghoniem, N. (2022).
\newblock A model of thermal creep and annealing in finite domains based on
coupled dislocation climb and vacancy diffusion.
\newblock \emph{Journal of the Mechanics and Physics of Solids}, \textbf{169}, 105066.

\bibitem{Po2014MoDELib}
Po, G. and Ghoniem, N. (2014).
\newblock A variational formulation of constrained dislocation dynamics coupled
with heat and vacancy diffusion (MoDELib --- Mechanics of Defect Evolution
Library).
\newblock \emph{Journal of the Mechanics and Physics of Solids}, \textbf{66}, 103--116.

\bibitem{Woo1988}
Woo, C. H. (1988).
\newblock Theory of irradiation deformation in non-cubic metals: effects of
anisotropic diffusion.
\newblock \emph{Journal of Nuclear Materials}, \textbf{159}, 237--256.

% ---- numerical multiple-time-scale references ----
\bibitem{Tikhonov1952}
Tikhonov, A. N. (1952).
\newblock Systems of differential equations containing small parameters in the
derivatives.
\newblock \emph{Matematicheskii Sbornik}, \textbf{31}(73), 575--586.

\bibitem{Fenichel1979}
Fenichel, N. (1979).
\newblock Geometric singular perturbation theory for ordinary differential
equations.
\newblock \emph{Journal of Differential Equations}, \textbf{31}, 53--98.

\bibitem{SegelSlemrod1989}
Segel, L. A. and Slemrod, M. (1989).
\newblock The quasi-steady-state assumption: a case study in perturbation.
\newblock \emph{SIAM Review}, \textbf{31}(3), 446--477.

\bibitem{MaasPope1992}
Maas, U. and Pope, S. B. (1992).
\newblock Simplifying chemical kinetics: intrinsic low-dimensional manifolds in
composition space.
\newblock \emph{Combustion and Flame}, \textbf{88}, 239--264.

\bibitem{LamGoussis1994}
Lam, S. H. and Goussis, D. A. (1994).
\newblock The CSP method for simplifying kinetics.
\newblock \emph{International Journal of Chemical Kinetics}, \textbf{26}, 461--486.

\bibitem{EEngquist2003}
E, W. and Engquist, B. (2003).
\newblock The heterogeneous multiscale methods.
\newblock \emph{Communications in Mathematical Sciences}, \textbf{1}(1), 87--132.

\bibitem{Kevrekidis2003}
Kevrekidis, I. G., Gear, C. W., Hyman, J. M., Kevrekidis, P. G., Runborg, O.,
and Theodoropoulos, C. (2003).
\newblock Equation-free, coarse-grained multiscale computation: enabling
microscopic simulators to perform system-level analysis.
\newblock \emph{Communications in Mathematical Sciences}, \textbf{1}(4), 715--762.

\bibitem{GearKevrekidis2003}
Gear, C. W. and Kevrekidis, I. G. (2003).
\newblock Projective methods for stiff differential equations: problems with
gaps in their eigenvalue spectrum.
\newblock \emph{SIAM Journal on Scientific Computing}, \textbf{24}(4), 1091--1106.

\bibitem{AscherRuuthSpiteri1997}
Ascher, U. M., Ruuth, S. J., and Spiteri, R. J. (1997).
\newblock Implicit-explicit Runge-Kutta methods for time-dependent partial
differential equations.
\newblock \emph{Applied Numerical Mathematics}, \textbf{25}, 151--167.

\bibitem{Strang1968}
Strang, G. (1968).
\newblock On the construction and comparison of difference schemes.
\newblock \emph{SIAM Journal on Numerical Analysis}, \textbf{5}(3), 506--517.

\bibitem{Chinesta2011}
Chinesta, F., Ammar, A., and Cueto, E. (2011).
\newblock Recent advances and new challenges in the use of the proper
generalized decomposition for solving multidimensional models.
\newblock \emph{Archives of Computational Methods in Engineering}, \textbf{18}, 395--404.

\end{thebibliography}

\end{document}
